% Vietnamese translation for OI Public Library
% Source-File: IOI中国国家候选队论文/2007/day1/10.陈启峰《Size Balanced Tree》.ppt
\documentclass[11pt]{article}
\usepackage{oipl-vn}

\newcommand{\Oh}{\mathcal{O}}

\title{Bản trình chiếu: Size Balanced Tree}
\author{Chen Qifeng}
\date{2007}

\begin{document}
\maketitle

\origtitle{Slide companion for Size Balanced Tree}
\sourcepath{IOI China candidate team papers / 2007 / day 1 / Chen Qifeng.ppt}

\section{Phạm vi bản dịch}

Tệp PowerPoint là bản trình chiếu của bài viết \textit{Size Balanced Tree} đã
được dịch từ PDF. Các slide dùng nhiều sơ đồ cây và mã giả ngắn; bản ghi chú
này tập trung vào thuật ngữ, bất biến và các thao tác cốt lõi.

\section{Cấu trúc dữ liệu}

Size Balanced Tree, viết tắt là SBT, là cây tìm kiếm nhị phân được giữ cân
bằng bằng kích thước cây con. Mỗi nút lưu khóa, con trái, con phải và
\(s[t]\), trong đó \(s[t]\) là số nút trong cây con gốc \(t\). Con trỏ rỗng
được xem có kích thước \(0\).

Hai thao tác cục bộ nền tảng là quay trái và quay phải. Phép quay giữ nguyên
thứ tự inorder của cây tìm kiếm nhị phân, nhưng thay đổi gốc cục bộ để khôi
phục quan hệ cân bằng theo kích thước.

\section{Bất biến cân bằng}

Với nút \(T\), gọi \(L\) và \(R\) lần lượt là con trái và con phải. Nếu
\(A,B\) là hai cây con của \(L\), còn \(C,D\) là hai cây con của \(R\), SBT
duy trì trực giác sau:
\[
  s[A],s[B]\le s[R],
  \qquad
  s[C],s[D]\le s[L].
\]
Khi một trong các bất đẳng thức bị phá vỡ, thủ tục \textsc{Maintain} chọn
trường hợp tương ứng và thực hiện một hoặc hai phép quay.

\section{Các thao tác}

Các slide liệt kê các thao tác chuẩn của tập động:
\textsc{Insert}, \textsc{Delete}, \textsc{Find}, \textsc{Rank},
\textsc{Select}, \textsc{Pred} và \textsc{Succ}. Do mỗi thao tác đi theo một
đường từ gốc rồi gọi \textsc{Maintain} ở các vị trí cần thiết, độ phức tạp được
giữ ở mức \(\Oh(\log n)\).

\section{Ghi chú triển khai}

Điểm mạnh của SBT trong thi lập trình là mã ngắn và dễ mở rộng. Trường
\(s[t]\) vừa phục vụ cân bằng, vừa cho phép trả lời truy vấn thứ hạng và chọn
phần tử thứ \(k\). Khi cài đặt, điều quan trọng là cập nhật kích thước đúng
ngay sau mỗi phép quay và sau mỗi lần chèn hoặc xóa.

\end{document}
